% GENERATED FILE. Edit canonical lessons, then run npm run generate:books.
% canonical-source-hash: fnv1a64:36ae464d0c7c6791
% canonical-lessons: GU-C190-khatu, GU-C190-jama-karvu, GU-C190-cek, GU-C190-vatavvu, GU-C190-udhar-levu

\chapter{At the bank counter}
\label{ch:gu-at-the-bank-counter}
\hlchaptermodality{190}

\begin{chapteropening}
\textbf{By the end of this chapter:} \emph{I can read bank words for accounts, deposits and cheques.}

Complete the last lesson of chapter 190: I can read bank words for accounts, deposits and cheques.
\end{chapteropening}

\section[\texorpdfstring{\gu{ખાતું} (khātũ)}{ખાતું (khātũ)}]{\gu{ખાતું} (khātũ) — bank account, account}

\label{lesson:GU-C190-khatu}



\begin{quote}
Before the new one: say the Gujarati for to be left over, then the Gujarati for to accept.
\end{quote}

\subsection*{You'll want to know: \gu{ખાતું}}
\textbf{\gu{ખાતું}} — \emph{khātũ} — \textquotedblleft{}bank account, account\textquotedblright{}.

\begin{quote}
\textbf{\gu{મારું} \gu{બેંકમાં} \gu{ખાતું} \gu{છે}.} — \emph{mārũ bẽkmā̃ khātũ che} — \textquotedblleft{}I have an account at the bank.\textquotedblright{}
\end{quote}

A neuter noun.

\begin{cousinweb}
\textbf{\gu{ખાતું}} (\emph{khātũ}) is a tadbhava. CDIAL lists it under Sanskrit \emph{kṣatrá}, \textquotedblleft{}might, rule\textquotedblright{}, with the Gujarati senses \textquotedblleft{}administrative department, account, account book\textquotedblright{} (CDIAL 3648).
\end{cousinweb}

\begin{grammarlens}[title={What you've built}]
One more word for everyday life at A2.
\end{grammarlens}

\subsection*{Your turn}
\begin{itemize}
\raggedright
  \item \emph{Say it:} \emph{khātũ}
  \item \emph{Say it:} \emph{khātũ}, once more
  \item \emph{Say it:} \emph{mārũ bẽkmā̃ khātũ che} — I have an account at the bank.
  \item \emph{Recall:} say the Gujarati for to be left over, then the Gujarati for to accept, then say \emph{khātũ} again
\end{itemize}

\begin{tcolorbox}[breakable,colback=teal!4,colframe=teal!35!black,title={Before you move on}]
What does \gu{ખાતું} mean? (\textquotedblleft{}Bank account, account\textquotedblright{}.) Say it once more.
\end{tcolorbox}
\section[\texorpdfstring{\gu{જમા} \gu{કરવું} (jamā karvũ)}{જમા કરવું (jamā karvũ)}]{\gu{જમા} \gu{કરવું} (jamā karvũ) — to deposit}

\label{lesson:GU-C190-jama-karvu}



\begin{quote}
Before the new one: say the Gujarati for to accept, then the Gujarati for bank account.
\end{quote}

\subsection*{You'll want to know: \gu{જમા} \gu{કરવું}}
\textbf{\gu{જમા} \gu{કરવું}} — \emph{jamā karvũ} — \textquotedblleft{}to deposit\textquotedblright{}.

\begin{quote}
\textbf{\gu{હું} \gu{બેંકમાં} \gu{પૈસા} \gu{જમા} \gu{કરું} \gu{છું}.} — \emph{hũ bẽkmā̃ paisā jamā karũ chũ} — \textquotedblleft{}I deposit money in the bank.\textquotedblright{}
\end{quote}

\begin{cousinweb}
\textbf{\gu{જમા} \gu{કરવું}} (\emph{jamā karvũ}) is built from \emph{jamā}, \textquotedblleft{}deposited, credited\textquotedblright{}, a loan from Arabic \emph{jamʿ}, \textquotedblleft{}gathering, sum\textquotedblright{}, that reached Gujarati through Persian, and the light verb \emph{karvũ}, \textquotedblleft{}to do\textquotedblright{}. Its Arabic root is \emph{j-m-ʿ}, \textquotedblleft{}to gather\textquotedblright{} (Lane).
\end{cousinweb}

\begin{grammarlens}[title={What you've built}]
One more word for everyday life at A2.
\end{grammarlens}

\subsection*{Your turn}
\begin{itemize}
\raggedright
  \item \emph{Say it:} \emph{jamā karvũ}
  \item \emph{Say it:} \emph{jamā karvũ}, once more
  \item \emph{Say it:} \emph{hũ bẽkmā̃ paisā jamā karũ chũ} — I deposit money in the bank.
  \item \emph{Recall:} say the Gujarati for to accept, then the Gujarati for bank account, then say \emph{jamā karvũ} again
\end{itemize}

\begin{tcolorbox}[breakable,colback=teal!4,colframe=teal!35!black,title={Before you move on}]
What does \gu{જમા} \gu{કરવું} mean? (\textquotedblleft{}To deposit\textquotedblright{}.) Say it once more.
\end{tcolorbox}
\section[\texorpdfstring{\gu{ચેક} (cek)}{ચેક (cek)}]{\gu{ચેક} (cek) — cheque}

\label{lesson:GU-C190-cek}



\begin{quote}
Before the new one: say the Gujarati for bank account, then the Gujarati for to deposit.
\end{quote}

\subsection*{You'll want to know: \gu{ચેક}}
\textbf{\gu{ચેક}} — \emph{cek} — \textquotedblleft{}cheque\textquotedblright{}.

\begin{quote}
\textbf{\gu{તેણે} \gu{મને} \gu{ચેક} \gu{આપ્યો}.} — \emph{teṇe mane cek āpyo} — \textquotedblleft{}He gave me a cheque.\textquotedblright{}
\end{quote}

A masculine noun.

\begin{cousinweb}
\textbf{\gu{ચેક}} (\emph{cek}) is a loan from English \emph{cheque}.
\end{cousinweb}

\begin{grammarlens}[title={What you've built}]
One more word for everyday life at A2.
\end{grammarlens}

\subsection*{Your turn}
\begin{itemize}
\raggedright
  \item \emph{Say it:} \emph{cek}
  \item \emph{Say it:} \emph{cek}, once more
  \item \emph{Say it:} \emph{teṇe mane cek āpyo} — He gave me a cheque.
  \item \emph{Recall:} say the Gujarati for bank account, then the Gujarati for to deposit, then say \emph{cek} again
\end{itemize}

\begin{tcolorbox}[breakable,colback=teal!4,colframe=teal!35!black,title={Before you move on}]
What does \gu{ચેક} mean? (\textquotedblleft{}Cheque\textquotedblright{}.) Say it once more.
\end{tcolorbox}
\section[\texorpdfstring{\gu{વટાવવું} (vaṭāvvũ)}{વટાવવું (vaṭāvvũ)}]{\gu{વટાવવું} (vaṭāvvũ) — to cash a cheque, to exchange}

\label{lesson:GU-C190-vatavvu}



\begin{quote}
Before the new one: say the Gujarati for to deposit, then the Gujarati for cheque.
\end{quote}

\subsection*{You'll want to know: \gu{વટાવવું}}
\textbf{\gu{વટાવવું}} — \emph{vaṭāvvũ} — \textquotedblleft{}to cash a cheque, to exchange\textquotedblright{}.

\begin{quote}
\textbf{\gu{તે} \gu{બેંકમાં} \gu{ચેક} \gu{વટાવે} \gu{છે}.} — \emph{te bẽkmā̃ cek vaṭāve che} — \textquotedblleft{}She cashes a cheque at the bank.\textquotedblright{}
\end{quote}

\begin{cousinweb}
\textbf{\gu{વટાવવું}} (\emph{vaṭāvvũ}) is a tadbhava. It comes from Sanskrit \emph{vartáyati}, \textquotedblleft{}causes to turn\textquotedblright{}, and CDIAL glosses the Gujarati verb \textquotedblleft{}to exchange, cash\textquotedblright{} (CDIAL 11356).
\end{cousinweb}

\begin{grammarlens}[title={What you've built}]
One more word for everyday life at A2.
\end{grammarlens}

\subsection*{Your turn}
\begin{itemize}
\raggedright
  \item \emph{Say it:} \emph{vaṭāvvũ}
  \item \emph{Say it:} \emph{vaṭāvvũ}, once more
  \item \emph{Say it:} \emph{te bẽkmā̃ cek vaṭāve che} — She cashes a cheque at the bank.
  \item \emph{Recall:} say the Gujarati for to deposit, then the Gujarati for cheque, then say \emph{vaṭāvvũ} again
\end{itemize}

\begin{tcolorbox}[breakable,colback=teal!4,colframe=teal!35!black,title={Before you move on}]
What does \gu{વટાવવું} mean? (\textquotedblleft{}To cash a cheque, to exchange\textquotedblright{}.) Say it once more.
\end{tcolorbox}
\section[\texorpdfstring{\gu{ઉધાર} \gu{લેવું} (udhār levũ)}{ઉધાર લેવું (udhār levũ)}]{\gu{ઉધાર} \gu{લેવું} (udhār levũ) — to borrow}

\label{lesson:GU-C190-udhar-levu}



\begin{quote}
Before the new one: say the Gujarati for cheque, then the Gujarati for to cash a cheque.
\end{quote}

\subsection*{You'll want to know: \gu{ઉધાર} \gu{લેવું}}
\textbf{\gu{ઉધાર} \gu{લેવું}} — \emph{udhār levũ} — \textquotedblleft{}to borrow\textquotedblright{}.

\begin{quote}
\textbf{\gu{તે} \gu{બેંકમાંથી} \gu{પૈસા} \gu{ઉધાર} \gu{લે} \gu{છે}.} — \emph{te bẽkmā̃thī paisā udhār le che} — \textquotedblleft{}He borrows money from the bank.\textquotedblright{}
\end{quote}

Past \textbf{\gu{તેણે} \gu{ઉધાર} \gu{લીધું}}, \emph{teṇe udhār līdhũ}.

\begin{cousinweb}
\textbf{\gu{ઉધાર} \gu{લેવું}} (\emph{udhār levũ}) is built from \emph{udhār}, \textquotedblleft{}credit, loan\textquotedblright{}, and \emph{levũ}, \textquotedblleft{}to take\textquotedblright{}. \emph{udhār} is a tadbhava from Sanskrit \emph{uddhārá}, \textquotedblleft{}raising; a loan\textquotedblright{}, and CDIAL lists Gujarati \emph{udhāro}, \textquotedblleft{}credit account\textquotedblright{} (CDIAL 2018).
\end{cousinweb}

\begin{grammarlens}[title={What you've built}]
One more word for everyday life at A2.
\end{grammarlens}

\subsection*{Your turn}
\begin{itemize}
\raggedright
  \item \emph{Say it:} \emph{udhār levũ}
  \item \emph{Say it:} \emph{udhār levũ}, once more
  \item \emph{Say it:} \emph{te bẽkmā̃thī paisā udhār le che} — He borrows money from the bank.
  \item \emph{Recall:} say the Gujarati for cheque, then the Gujarati for to cash a cheque, then say \emph{udhār levũ} again
\end{itemize}

\begin{tcolorbox}[breakable,colback=teal!4,colframe=teal!35!black,title={Before you move on}]
What does \gu{ઉધાર} \gu{લેવું} mean? (\textquotedblleft{}To borrow\textquotedblright{}.) Say it once more.
\end{tcolorbox}
